%% ma: L11-B15-0031
%% lop: 11
%% chuong: 5
%% bai: 15
%% dang: 11.15.4
%% loai: TN4
%% muc: TH
%% dap_an: "B"
%% nguon: "tai-lieu-hoc-tap-mon-toan-lop-11-theo-chuong-trinh-giao-duc-pho-thong-2018.pdf (Lớp Toán Thầy Duy, Phạm Lê Duy, THPT Chu Văn An), Chương 3 Bài 1, Phần 1, Câu 20 (tr.322)"
%% kien_thuc: "Lớp 11 Bài 15 (giới hạn của dãy số)"
%% trang_thai: chinh-thuc
%% ghi_chu: "Đã duyệt 03/11/2026."
\begin{ex}
$\lim\dfrac{100^{n+1}+3\cdot99^n}{10^{2n}-2\cdot98^{n+1}}$ là
\choice{$+\infty$}{\True $100$}{$\dfrac1{100}$}{$0$}
\loigiai{\sloppy Vì $10^{2n}=100^n$, chia tử và mẫu cho $100^n$: \[\dfrac{100+3\left(\frac{99}{100}\right)^n}{1-196\left(\frac{49}{50}\right)^n}\to\dfrac{100}{1}=100\] (do $2\cdot98^{n+1}=196\cdot98^n$ và $\frac{98}{100}=\frac{49}{50}$).}
\end{ex}
